\section{Introduction}\label{sec:introduction}

When a compact Riemann surface develops nodes, a metric
on a bundle can equilibrate on the normalization components
long before the components equilibrate with one another.
The remaining variables are matrix metrics. Their coupling
depends on the fibre identifications at the nodes, including
cycle holonomy that cannot be removed by component frame
changes. We prove a finite nonlinear reduction of the
actual Hermitian--Yang--Mills metric flow in this setting.
The comparison allows independent node lengths and
noncommuting unitary gluing. It is accompanied by a complete
relative spectral theorem, entry from spatial initial
metrics, and geometric and inverse consequences.

The principal family smooths four projective lines in
general position. Its six nodes have independent exact
holomorphic equations $u_ev_e=\tau_e$; its smooth fibres
$C_\tau$ are plane quartics of genus three. We use the
metric induced by the Fubini--Study form, normalized so
that a line has area one. The construction in
Section~\ref{sec:geometry} gives a proper projective
family and, for every tuple $U=(U_e)\in\SU(r)^6$,
a holomorphic bundle with flat unitary connection $D_U$,
parallel metric $h_0$ and trivialized determinant.
Here $r\ge2$ is fixed. The component connections are
trivial, and fibre holonomy factors through the free
group of rank three of the dual graph $K_4$.
The bundles are polystable of degree zero and are
stable exactly when this representation is irreducible.

\subsection{The effective equation and the main theorem}
Transport $U_{ij}$ takes the $j$ component frame to the
$i$ frame. Set
\[
 L_e=\log(1/|\tau_e|),\qquad L_*=\min_eL_e,\qquad
 w_e=2\pi/L_e .
\]
All parameters are fixed during each flow, and the
degeneration limit is $L_*\to\infty$.
Let $A_i$ denote the limiting component areas, equal
to one for the chosen Fubini--Study metric. On the
real space of trace-free Hermitian tuples
$\cV=\bigoplus_{i=0}^3\Herm_0(r)$ define
\[
 M(X)=\sum_i A_i\tr(X_i^2),\qquad
 Q_{\tau,U}(X)=\sum_{e=ij}w_e
       \|X_i-U_{ij}X_jU_{ij}^*\|_{\HS}^2.
\]
Let $\nu_0\le\cdots\le\nu_{m-1}$ be the generalized
eigenvalues of $(Q_{\tau,U},M)$, including zeros,
where $m=4(r^2-1)$. The actual linearized operator is
$\Delta_{\tau,U}=D_U^*D_U$ on trace-free
$h_0$-Hermitian endomorphisms, with the unnormalized
area norm. Its eigenvalues are denoted by $\lambda_k$.

Write $\cP_r$ for the determinant-one positive
Hermitian matrices, with symmetric-space metric
$g_P(V,W)=\tr(P^{-1}VP^{-1}W)$ and distance $d$.
The finite potential and its negative gradient
equation, using the area-weighted product metric, are
\begin{equation}\label{eq:intro-finite}
 \begin{split}
 \Phi_{\tau,U}(\mathbf H)
   &=\frac12\sum_{e=ij}w_e
                d(H_i,U_{ij}H_jU_{ij}^*)^2,\\
 H_i^{-1}\dot H_i
   &=\frac1{A_i}\sum_{e=ij}w_e
                  \log(H_i^{-1}U_{ij}H_jU_{ij}^*).
 \end{split}
\end{equation}
The exact metric flow uses physical time $s$ and
the normalization $H^{-1}\partial_sH=-2i\Lambda F_H$.
For two metrics set
$\sigma(H,K)=\tr(H^{-1}K+K^{-1}H)-2r$.
It is intrinsic under frame changes and bounds
$d(H,K)^2$ from above.

We also specify an initial class that does not
require an already matched finite tuple. Put
$\rho=L_*^{-1}$, fix $0<\beta<1$ and $R_0>0$,
and write an initial metric as $h_0e^S$.
Let $\Omega_i(2\rho)\subset C_i$ be the complement
of the incidence discs of radius $2\rho$.
Pull $S$ back to these component exteriors using
the fixed geometric identifications and flat
frames of Section~\ref{sec:geometry}. Require
\begin{equation}\label{eq:intro-initial-class}
 \|S\|_{C^0(C_\tau)}
 +\inf_{\widetilde S_i|_{\Omega_i(2\rho)}=S_i}
           \sum_i\|\widetilde S_i\|_{C^{6,\beta}(C_i)}
 <R_0 ,
\end{equation}
where the extensions are trace-free Hermitian
matrices on the fixed components.
This is a nonempty open class on each smooth
fibre. Its radius is independent of the degeneration;
bounded smooth data on the middle necks may have
arbitrarily large derivatives.

A unitary frame identifies each endomorphism fibre
with $\operatorname{Mat}_r(\CC)$ and so specifies
the matrix-algebra marking used below.
Projective transport means the class of a unitary
edge map in $\operatorname{PSU}(r)=\SU(r)/Z(\SU(r))$.
After choosing a base vertex, cycle transports
are represented in its frame; changing that
frame conjugates all cycle maps simultaneously.

\begin{theorem}\label{thm:main}
The compact quartic family and flat bundles above
have the following properties.
\begin{enumerate}
\renewcommand{\labelenumi}{(\roman{enumi})}
\item There are $\gamma>0$ and $\epsilon(\tau)\to0$
as $L_*\to\infty$ such that
\[
 (1-\epsilon(\tau))\nu_k\le\lambda_k
 \le(1+\epsilon(\tau))\nu_k\quad(0\le k<m),
 \qquad \lambda_m\ge\gamma .
\]
The estimates are uniform in all node-length ratios
and all $U$. The kernel is naturally identified with the space of
tuples satisfying $X_i=\Ad(U_{ij})X_j$.

\item For finite initial tuples in a fixed bounded
ball, \eqref{eq:intro-finite} has a confined
all-time solution. It determines a smooth
determinant-one comparison metric $\widehat K(s)$
on the entire $C_\tau$, converging to $H_i(s)$
on fixed component compact subsets with error
$O(L_*^{-1})$. If the exact initial metric has
fixed positive bounds against $h_0$ and
$\eta=\int_{C_\tau}\sigma(H^{\mathrm{in}},\widehat K(0))\,dA_\tau$,
then, with $e_*=(1+\log L_*)^2/L_*$,
\[
 \sup_{0\le s\le TL_*}\int_{C_\tau}
       d(H(s),\widehat K(s))^2\,dA_\tau
       \le\eta+CTe_*\qquad(T\ge0).
\]
The constants are independent of $T$, the
node-length ratios and the unitary transports.

\item For the spatial initial class
\eqref{eq:intro-initial-class}, the component
flows select a bounded tuple $P$ and there is
a uniform preparation time $s_*\to\infty$,
$s_*/L_*\to0$, such that the comparison in
part~{\rm(ii)} applies to $H(s_*+s)$ and the
finite flow starting at $P$, with initial
error $\eta_*\to0$ uniformly over that class.

\item If $L_*\nu_0\ge\kappa>0$, then for the
initial data and discrepancy in part~{\rm(ii)}
\[
 \sup_{s\ge0}\int_{C_\tau}d(H(s),\widehat K(s))^2\,dA_\tau
 \le C\eta+C_\kappa e_*(1+|\log e_*|).
\]
The same conclusion holds after the preparation
in part~{\rm(iii)}, with $\eta_*$ in place of $\eta$.

\item On the common-node slice $\tau_e=t$, put
$L=\log(1/|t|)$. The graph representation
quotient $\SU(r)^{6,\mathrm{irr}}/\SU(r)^4$
embeds as a Lagrangian locus in the irreducible
flat-bundle moduli space of $C_t$.
Its pullback Hodge squared norm is $2\pi/L$
times the quotient of the bi-invariant product
squared norm on $\SU(r)^6$, with relative
error $O(L^{-1})$ uniform towards reducibility.
Moreover, the first $m$ eigenvalues and the
component averages of their $L^2$-orthonormal
eigensections, in fixed unitary matrix-algebra
frames, consistently determine the projective
edge transports. Cycle transports are recovered
up to root conjugation. Central lifts are not
determined by these observations.
\end{enumerate}
\end{theorem}
\begin{proof}
Propositions~\ref{prop:compact-family} and
\ref{prop:bundle} construct the family, bundle
and holonomy representation. Theorem~\ref{thm:spectrum}
proves part~(i).
Proposition~\ref{prop:finite-flow} and
Theorem~\ref{thm:tracking} prove part~(ii);
Theorems~\ref{thm:component-phase} and
\ref{thm:entry} give the component limits and preparation
in part~(iii).
Theorem~\ref{thm:all-time} gives part~(iv).
The moduli statement is Theorem~\ref{thm:moduli-metric},
and the observed-operator construction of
Lemma~\ref{lem:marked-operator} together with
Theorem~\ref{thm:projective-inverse} gives the
last assertion. All these constructions use
the same quartic family, flat bundles and
metric normalization.
\end{proof}

The relative estimate in part~(i) controls every
positive small eigenvalue even when its ratio
to the largest weight tends to zero. It does
not imply nonlinear tracking on the inverse
of that weakest eigenvalue. For example,
Proposition~\ref{prop:three-scales} realizes
orders $h,h^2,h^4$ on actual stable rank-two
bundles. The uniform all-time hypothesis in
part~(iv) fails in that example.
Likewise, part~(v) concerns the indicated
moduli locus and marked observations. It does
not assert an ambient moduli completion or
reconstruction from eigenvalues alone.

\subsection{Component holonomy and integral monodromy}
Two further constructions test different aspects
of the reduction. In Section~\ref{sec:component-holonomy}
a compact genus-three family has two elliptic
and two rational normalization components.
The component flat connections vary with
parameters $\alpha_i$ independently of the
common node parameter. The effective energy
contains the component commutator terms of
order $\alpha_i^2$ and the transported edge
terms of order $L^{-1}$.
Theorems~\ref{thm:coupled-spectrum},
\ref{thm:coupled-tracking} and
\ref{thm:coupled-entry} prove the complete
relative spectrum, actual nonlinear comparison
and spatial entry for arbitrary ratios
of these scales. The nonlinear error on
bounded intervals in $s/L$ tends to zero
as $\max_i|\alpha_i|+L^{-1}\to0$.
This family has a common node scale;
it does not combine varying component
holonomy with all six independent quartic
parameters.

The topology problem in Section~\ref{sec:inverse}
starts with six labelled rational components,
each with three unpaired branch points.
Their pairing is unknown. We construct
compact first-order smoothings for a finite
class of simple trivalent pairings, with
identical prescribed component metrics and
local node metric germs. The first
$6(r^2-1)$ eigenvalues and component averages
of the eigensections determine the missing
adjacency for all sufficiently small parameters.
The known first-order smoothing and positive
base-loop orientation then determine the
integral symplectic conjugacy class of
monodromy on first homology.

For the triangular prism and $K_{3,3}$ the
torsion groups of the cokernel of monodromy minus the identity are
$\mathbb Z/5\mathbb Z\oplus\mathbb Z/15\mathbb Z$
and $(\mathbb Z/3\mathbb Z)^2\oplus\mathbb Z/9\mathbb Z$.
Both smooth fibres have genus four; the
distinguished invariant belongs to the
degeneration, not to their diffeomorphism
types. The classical Picard--Lefschetz and
graph-lattice identifications
\cite{achinger-ogus,bacher-harpe-nagnibeda}
are used after the analytic observations
have supplied the missing incidence.
These same compact families also carry
the noncommuting HYM dynamics, by
Proposition~\ref{prop:fixed-graph-extension}.

\subsection{Relation to earlier work and the proof}
Relative graph descriptions of small surface
eigenvalues have a substantial history. Burger
\cite[Theorem~1.1]{burger} proves relative scalar
asymptotics on degenerating hyperbolic surfaces
without a lower bound on the ratios of the short
geodesic lengths. Thus neither a graph energy nor
arbitrarily separated scalar scales is new by itself.
Our surface metrics are induced from a smooth ambient
metric, the coefficients act on bundle endomorphisms,
and the holonomy kernel may change. Although the
Dirichlet energy in two dimensions is conformally
invariant, the function $L^2$ mass is not; the
ambient-metric normalization requires its own
comparison.

For the spectral theorem, component averages
give the finite variables and logarithmic
interpolation gives trial sections with the
prescribed edge energy. The reverse comparison
bounds each boundary mismatch by the energy
of the actual section before summing with the
individual edge weights. This avoids division
by the smallest weight. A uniform annular
mass estimate rules out additional weak modes;
min--max then compares the entire small
cluster relatively, including changing
zero eigenspaces.

Daskalopoulos--Wentworth \cite{dw} study connection
spectra and Yang--Mills flow near the boundary of
Teichm\"uller space. The component-trivial connections
considered here have a large limiting parallel
space even when the global representation is
irreducible. We estimate the slow coupling of that
space and its nonlinear metric evolution.
The fixed-base asymptotic theories of
Haiden--Katzarkov--Kontsevich--Pandit \cite{hkkp}
and Lam \cite{lam} provide a complementary
reduction of HYM evolution. Their fixed-background
estimates do not alone give uniform constants
as the underlying curve pinches.

The free-group construction of the moduli locus is
the classical unitary Schottky construction
\cite{florentino}. Daskalopoulos--Wentworth
\cite{dw-local} analyze local degeneration of
bundle moduli by weighted elliptic methods.
Section~\ref{sec:moduli} explains the hypotheses
of the relevant kernel-continuation statements
and proves the present relative Hodge-norm
comparison directly by minimizing the energy
of affine primitives. The uniform estimate does
not require a lower bound for a nonzero
eigenvalue of the vertex-gauge operator.
Conformal changes of the surface metric alone
do not distinguish the one-form energy.

For trivial bundles and a common neck scale,
the actual noncommuting matrix-flow comparison
was established in \cite{wuebben-hym}.
We retain its full Chern-curvature equation and
intrinsic trace comparison, and prove the
weighted gluing estimates needed for independent
node lengths and nontrivial unitary transport.
The regular component correction at a node
receives contributions from every incident
edge, whereas the node profile has its own
logarithmic slope. Keeping both weights in
the collar estimate gives the uniform residual
$O(L_*^{-2}(1+\log L_*)^2)$.
The trace comparison integrates this residual
without a degeneration-dependent Sobolev
constant or exponential time factor.
For the all-time estimate we combine this
finite-window comparison with exponential
decay of both the exact and finite solutions.
The lower bound for $L_*\nu_0$ makes these
tails uniform in the rescaled time.
The finite gradient structure does not identify
the actual HYM equation with ordinary
harmonic-map heat flow; its additional
commutator term is retained in
\eqref{eq:real-hym}.

Finally, the spatial entry argument uses a
fixed-component entropy: the logarithm of
the determinant of the arithmetic matrix mean.
It selects the constant limit metric with an
exponential estimate, after which a global
preparation metric transfers the fixed-component
flows to the smoothing. Convergence at a
fixed positive rescaled time also requires
convergence of these component limits and
the normalized finite parameters.
Uniform tracking alone does not impose it.

For the family with varying component
holonomy, a harmonic gauge on each punctured
component separates its mean from a controlled
remainder. The boundary error is bounded by
the component energy, rather than by the
size of its nearly parallel mean. Summing
the disjoint component and neck energies
then preserves arbitrary ratios of their
scales. In the nonlinear comparison the
curvature of a spatially constant component
metric is canceled exactly by its finite
component velocity. The additional matching
error caused by component holonomy has order
$\max_i|\alpha_i|/L$, which
is small after integration on bounded
intervals in $s/L$.

The inverse argument first uses component
averages to identify the actual low spectral
subspace with the vertex space. Polar
normalization makes that identification
isometric. The lower energy estimate gives
a positive operator difference, and relative
eigenvalue comparison bounds its trace;
together these control the observed operator.
On a known simple graph its off-diagonal
blocks recover the adjoint transports.
When adjacency is unknown, the block norms
distinguish edges from nonedges. The recovered
graph and the prescribed smoothing orders
then supply the integral cycle pairing
in the classical monodromy formula.
